\section{Testing of the New Implementation}
\label{sec:testing}

\subsection{Test Strategy}
\todo{Unit vs.\ integration tests, use of test doubles (fake I/O, seeded
\code{Random}, injectable deck), how the design makes this possible.}

\subsection{Test Cases per Requirement}
\begin{longtable}{@{}p{1.2cm}p{5.5cm}p{5cm}p{2cm}@{}}
\toprule
\textbf{Req.} & \textbf{Test (class\#method)} & \textbf{What is verified} & \textbf{Result} \\
\midrule
\endhead
\req{1} & \todo{\code{GameSetupTest\#rejectsTooManyPlayers}} & \todo{} & \todo{Pass} \\
\req{2} & \todo{} & \todo{} & \todo{} \\
\req{3} & \todo{} & \todo{} & \todo{} \\
\bottomrule
\caption{Traceability between requirements and unit tests.}
\label{tab:test-traceability}
\end{longtable}

\subsection{Example Test}
\begin{lstlisting}[caption={Example JUnit test.},label={lst:test-example}]
@Test
void defuseReinsertsKittenAtChosenPosition() {
    // TODO
}
\end{lstlisting}

\subsection{Results and Coverage}
\todo{Summary of test results and, if available, coverage numbers.}
